\section{Sizing the buffer: a negative result}
\label{sec:sizing}

Before arriving at a cap and a discard we built the thing the problem statement
suggests: a controller that sizes each buffer to what its traffic needs. It
works, in the sense that it reaches parity with the better fixed buffer on both
workloads without being told which workload it is running. It is nonetheless not
the right mechanism, and the reason is worth more than the mechanism would have
been.

\subsection{The construction}

TCP's receive-side rule is $\rcvbuf = 2 \times (\text{bytes received in one
RTT})$: the RTT delimits one turn of the consumer, and the factor of two leaves
room for the next turn's arrivals while the current one is consumed. UDP has no
RTT, which is the usual reason given for the absence of UDP autotuning.

The RTT's role, though, is only to delimit a period over which demand can be
measured, and the receive queue supplies one: the interval from its rising above
half full to its falling below an eighth. Over that interval the queue's peak is
what arrived while the consumer worked through what it already had---which, when
the NIC is batching, is one moderation window's delivery. Sizing to twice it
reproduces TCP's rule with a clock UDP can read, and lands on
$\rxframes \times \texttt{pkt}$ without consulting the driver.

Two properties fall out. A socket whose consumer cannot keep up never sees its
queue empty, so it never closes an interval, never produces a target, and never
grows---the case where more buffer cannot help excludes itself, with no separate
test. And a cycle must be closed by \emph{both} halves: detecting only the drain
is insufficient, because the enqueue path runs on every arrival and a socket
whose queue is already empty would close a cycle on each packet, with a peak of
one packet, overwriting a correct target with nothing.

\subsection{Why it is not the answer}

Table~\ref{tab:nostatic} places it against the alternatives. It matches the
better fixed buffer on both workloads. It does not beat a fixed 1\,MB buffer
combined with the discard of \S\ref{sec:shed}, which matches it below the
receiving core's saturation point and exceeds it by 4--6\% above
(Figure~\ref{fig:shedmarg}).

The reason is the measurement that motivated the sizer in the first place. Past
the point where arrivals outpace the consumer, the queue is full whatever its
capacity, and the extra capacity only lengthens the time data waits before
copy-out---by which time, per \S\ref{sec:workingset}, it has been evicted.
Sizing up is the wrong response to overload, and an adaptive sizer is a machine
for sizing up.

There is a second, sharper way to see this. The demand figure is the queue's
peak occupancy, and occupancy is bounded by \rcvbuf. A larger buffer therefore
permits a deeper peak, which raises the estimate, which grows the buffer again;
the loop is stopped by the allowance rather than by need, and we measured it
reaching 9.4\,MB on a socket that performed better at 1. TCP does not have this
problem because \texttt{tcp\_rcv\_space\_adjust()} counts bytes \emph{copied to
the application}, which buffer size cannot inflate. Counting bytes released to
the consumer over the same interval would be the faithful translation and would
remove the inflation.

We did not pursue it, because by then the conclusion no longer depended on the
estimator. Even a perfectly calibrated sizer converges on ``as much buffer as
this socket can use'', and the measurement says the quantity worth minimising
is how much data sits between the producer and the consumer. The right target is
small, and a mechanism whose whole purpose is to find how large it may safely be
is aimed the wrong way.

\subsection{What the exercise established}

Two things, both of which the final design rests on.

That \emph{discarding cannot substitute for a cap.} Eight sockets at 8\,MB reach
22.7\,Gbps, and 29.4 with discarding enabled---an improvement, and still half of
the 55.8 the same eight reach at 1\,MB, still losing 47\% of what was offered.
Removing work does not remove capacity. Whatever else is done, something has to
bound what the buffers sum to.

And that \emph{a cap is all the sizing that is needed.} Between the cap and the
discard, the remaining freedom---where within the bound a given socket sits---is
worth at most the 0.2\,Gbps by which the sizer leads the fixed buffer below the
ceiling, and costs 4--6\% above it.
